\documentclass[11pt]{achemso}
\usepackage{geometry}                % See geometry.pdf to learn the layout options. There are lots.
\geometry{letterpaper}                   % ... or a4paper or a5paper or ... 
\usepackage{graphicx}
\usepackage{amssymb}
\usepackage{epstopdf}
\usepackage{amsmath}
\usepackage{cancel}
\usepackage{float}

\DeclareGraphicsRule{.tif}{png}{.png}{`convert #1 `dirname #1`/`basename #1 .tif`.png}

\title{Stark Shifts in Simulated Excited State Vibrational Spectra}
\author{C. Houferak}
\email{houferak@uw.edu}
\author{E. Rabe}
\email{rabee@uw.edu}
\author{J. J. Radler}
\email{jjradler@uw.edu}
\affiliation[University of Washington]{Department of Chemistry, University of Washington, Seattle, WA, USA}
\date{\today}                                           % Activate to display a given date or no date

\begin{document}
\maketitle
% I THINK WE NEED TO REWORK THE ABSTRACT TO REFLECT THE PROJECT PLAN WE 
% PRESENTED IN OUR TALK -- JJR 6/2 
\begin{abstract}
Our chosen system of study is the \emph{2-[(7-4-[N,N-Bis(4-methylphenyl)amino]phenyl-2,1,3-
benzothiadiazol-4-yl)methylene]propanedinitrile} (DTDCPB) chromophore prevalent in organic
photovoltaic materials research. DTDCPB has nitrile groups that make it ideal for
IR studies due to its strong signature stretching frequency. For the first component of our
project, we hope to simulate some ground state IR spectra of DTDCPB in the presence
of varying external fields. This can then be used to calculate the Stark shift and build a
calibration curve for the effect of external field on the Stark shift of this molecule. Furthermore,
we plan to probe the excited state of DTDCPB as well and see if we can look at Stark
effects in the excited state as well. The third step we hope to investigate is a pump-probe
transient IR study of this molecule where the excited state dynamics of DTDCPB can be
observed. Hopefully our understanding of Stark effects in the ground and excited states
will elucidate some of the dynamics of the transient IR study. \textbf{CURRENTLY 175 WORDS}
\end{abstract}

%---------------------------------------------------------------------------------------------------------------------------------

\section{Introduction}

As research in solar energy becomes increasingly important for environmental applications,
 several new strategies are being investigated to increase the efficiency of these devices. 
 Specificlly in organic photovoltaics (OPVs), one such strategy has been to use small molecules with 
 large ground-state dipole moments. It is believed that these dipoles can allow for enhanced hole 
 hopping and can minimize charge recombination (ref Griffith). One particularly promising molecule 
 is  2-[(7-4-[N,N-Bis(4-methylphenyl)amino]phenyl-2,1,3-benzothiadiazol-4-yl)methylene]
 propanedinitrile (DTDCPB) and has resulted in OPVs with power conversion effiencies of up to 9.6\%.  
 The nitrile groups present in the molecule can act as reporter moieties for vibrational spectroscopy. 
 By modeling the effect of an applied electric field on absorption spectrum (Stark Effect), it is then 
 possible to gather information regarding the local electric field during charge transfer events in a 
 photovoltaic device via excited state IR spectroscopy. The first requires a thorough understanding 
 of the electric field dependence of the ground state IR spectrum before probing the excited state.  \textbf{CURRENTLY 168 WORDS}

%---------------------------------------------------------------------------------------------------------------------------------
\section{Theory}

\par In order to obtain vibrational spectra of the full DTDCPB dimer system and model the experiment efficiently, a technique called Born-Oppenheimer Molecular Dynamics (BOMD) was applied first to a cyanide molecule (as a test case), subsequently to DTDCPB monomer, and finally to the DTDCPB dimer system. BOMD is an \emph{ab initio} molecular dynamics method which couples quantum mechanical calculation of the electronic structure and potentials to a classical treatment of nuclear dynamics. We chose to use BOMD for the larger molecular systems specifically because it is able to capture more of the physics of the system of interest more efficiently than the traditional method of calculating the Hessian to obtain the frequency eigenvalues. Specifically, BOMD is able to capture the anharmonicity in nuclear motion, which may only be included in purely \emph{ab initio} calculations by including the third derivative (third order correction) term in the Taylor expansion expression of the energy as a function of nuclear position and creating an anharmonicity matrix from the system of equations describing the potential energy surface (PES). This is, of course, more computationally expensive than only calculating a first time-derivative of the dipole moment as in BOMD.  \textbf{CURRENTLY 195 WORDS}

\subsection{Standard Frequency Calculations}

\par First, a brief description of how the \emph{ab initio} frequency calculations are performed in the \textbf{g09} distribution of Gaussian.  

\begin{equation}
\begin{split}
E(Q) = E_{0}(Q) + 
\label{equation:Taylor}
\end{split}
\end{equation}

\subsection{Born-Oppenheimer MD (BOMD)}

\par Now some nuts and bolts for how BOMD is performed (Summarize, but with some equations for specifics like in the presentation.)

\par The nuclear dynamics calculations needed to simulate vibrational spectra were carried out in various dipole field strengths defined in \emph{Gaussian} in order to ``calibrate" the Stark shifts in the product spectra. First the molecular geometry was optimized in \emph{Gaussian} at the \textit{B3LYP/6-311G(d,p)} level of theory in the presence of a static electric field generated by a dipole oriented along the principal $z$-axis. Then, thermal sampling is carried out in the same electric field to generate trajectories for nuclear dynamics in a separate job using the optimized geometry.


\subsection{The Stark Effect}

Some discussion and math about the Stark Shift will go here. 

\subsection{Obtaining Spectra from Dipole Orientation Time Series}

How to calculate spectra from dipole autocorrelation functions vs. how it's done in the standard frequency calculations (refer to the presentation, but find original references for this...) 


%---------------------------------------------------------------------------------------------------------------------------------

\section{Methodology}



\textbf{I THINK WE SHOULD PUT THE NUTS AND BOLTS IN THE THEORY SECTION AND JUST BLAZE THROUGH METHODOLOGY REFERRING TO WHAT WAS EXPLAINED IN THE THEORY SECTION AND SUBSECTIONS}


\subsection{Calibrating Stark Shift}

The Stark effect comes in when a field is applied. \\

Applied fields to observe Stark shifts\\

Calibration curve:\\
\begin{center}
	Apply external E field \\
	Measure IR spectrum\\
	Calibrate Stark shifts for use in Pump Probe\\
	Watch Charge transfer (local E field) \\
\end{center}

Easy to put in an external field (molecule comes into Equilibrium in the field) to see the IR \\

The molecule is geometry optimized, then frequency calculated\\

Static = x,y,z then value of the field\\

Geometry optimization in a static E field (google the warnings to explain it) \\

Include point charges (dummy charges)\\ 

Dipole along z-axis, put two point charges on the axis and vary the distance to vary magnitude of the field. \\


\subsection{Start at ground state Calibration Curve thingy for Stark Shifts Static E fields}
a.	Opt+freq $\longrightarrow$ postprocess for IR

\subsection{Excited state opt+freq $\longrightarrow$ postprocess for IR}
a.	Same process

\subsection{Model pump probe using Stark shift Calibration (dynamics, no ext field)}
a.	BOMD  $\longrightarrow$ freq  postprocess for IR
\begin{center}
??????
PROFIT!!!! 
\end{center}

%---------------------------------------------------------------------------------------------------------------------------------

\section{Computational Results}


%---------------------------------------------------------------------------------------------------------------------------------

\section{Discussion}


%---------------------------------------------------------------------------------------------------------------------------------

\section{Conclusion}


%---------------------------------------------------------------------------------------------------------------------------------

\end{document}  
